\section{Results}
\label{sec:results}

We evaluate 8 zero-shot foundation models against 10 trained and statistical baselines on identical windows:
Chronos-2 \citep{ansari2026chronos2} and Chronos-Bolt \citep{ansari2024chronos}, TimesFM 2.5
\citep{das2024timesfm}, Toto-2.0 \citep{cohen2025toto}, TiRex \citep{auer2025tirex}, Moirai-2
\citep{liu2026moirai2,woo2024moirai}, Granite TTM r2 \citep{ekambaram2024ttm} and Sundial
\citep{liu2025sundial}. All models see a
1{,}024-day context except TTM, whose daily variants cap at 512.

\subsection{Foundation models lead, but most margins are not separable}

Averaged over the 16 daily tasks, six of the eight foundation models outrank every trained baseline. Chronos-2 has
the best mean rank at leads 7 and 30, 2.77 of 18 models overall; TiRex is best at lead 1. The strongest trained
baseline is PatchTST, at 6.96 (Table~\ref{tab:leaderboard}). Friedman rejects equality decisively at every lead
($\chi^2 = 220.1$, $215.5$, $209.8$; $p < 10^{-34}$).

That ranking, however, does not license an ordering. The Nemenyi critical difference is 6.58. At leads 7 and 30, the six
leading foundation models and PatchTST all fall within one critical difference of the best
(Figure~\ref{fig:cd}). The group separates clearly from the weaker baselines; positions within it do not.

Pairwise, the picture is more sober still. We compared the best foundation model with the best baseline for each
task and lead, correcting across all 48 comparisons. The foundation model is significantly better in 14, worse in
5, and statistically indistinguishable in 29. Those wins concentrate at a one-day lead and on the smooth reanalysis
tracks. Partitioned by track type, foundation models take 11 of 24 comparisons on reanalysis and only 3 of 24 on
the stations themselves, all three of those at lead 1. Of 432 model-task-lead skill intervals, 272 exclude
climatology.

\begin{table}[htbp]
  \centering
  \footnotesize
  \setlength{\tabcolsep}{4.5pt}
  \caption{Main leaderboard over the 16 daily tasks. Ranks are Friedman average ranks across the tasks at each lead (1 is best, 18 models); \emph{Overall} is the mean of the three. CRPSS is mean skill against the same track's day-of-year climatology at a 7-day lead, so a positive value beats climatology and zero is climatology, averaged over the 8 station-observation tasks and the 8 reanalysis tasks separately. The Nemenyi critical difference is 6.58, so rank gaps smaller than that are not separable. Parameter counts apply to the zero-shot models; the baselines are fitted per series. TTM r2 is point-only and is scored on point skill.}
  \label{tab:leaderboard}
  \begin{tabular}{lrrrrrrr}
    \toprule
    & Params & \multicolumn{4}{c}{Average rank (1 = best)} & \multicolumn{2}{c}{Mean CRPSS, lead 7} \\
    \cmidrule(lr){3-6} \cmidrule(lr){7-8}
    Model & (M) & Lead 1 & Lead 7 & Lead 30 & Overall & Station & Reanalysis \\
    \midrule
    \multicolumn{8}{l}{\emph{Zero-shot foundation models}} \\
    Chronos-2 & \num{119.5} & \num{4.19} & \num{2.12} & \num{2.00} & \num{2.77} & \num{+0.035} & \num{+0.128} \\
    TiRex & \num{35.3} & \num{1.62} & \num{3.44} & \num{3.62} & \num{2.90} & \num{+0.025} & \num{+0.121} \\
    TimesFM 2.5 & \num{231.3} & \num{4.38} & \num{4.44} & \num{3.75} & \num{4.19} & \num{+0.024} & \num{+0.121} \\
    Toto-2.0 & \num{312.7} & \num{3.75} & \num{5.81} & \num{6.56} & \num{5.38} & \num{-0.015} & \num{+0.117} \\
    Moirai-2 & \num{11.4} & \num{4.69} & \num{6.31} & \num{7.06} & \num{6.02} & \num{+0.003} & \num{+0.110} \\
    Chronos-Bolt & \num{205.3} & \num{5.81} & \num{5.75} & \num{6.62} & \num{6.06} & \num{+0.016} & \num{+0.108} \\
    Sundial & \num{128.3} & \num{11.06} & \num{10.69} & \num{10.31} & \num{10.69} & \num{-0.088} & \num{+0.031} \\
    TTM r2 & \num{0.8} & \num{16.38} & \num{15.94} & \num{15.25} & \num{15.85} & \num{-0.374} & \num{-0.246} \\
    \addlinespace
    \multicolumn{8}{l}{\emph{Trained and statistical baselines}} \\
    PatchTST & n/a & \num{7.75} & \num{6.31} & \num{6.81} & \num{6.96} & \num{+0.009} & \num{+0.080} \\
    AutoARIMA & n/a & \num{7.62} & \num{8.62} & \num{8.88} & \num{8.38} & \num{-0.062} & \num{+0.025} \\
    NHITS & n/a & \num{8.12} & \num{8.94} & \num{8.88} & \num{8.65} & \num{-0.002} & \num{+0.023} \\
    LightGBM & n/a & \num{12.88} & \num{8.62} & \num{8.25} & \num{9.92} & \num{-0.063} & \num{+0.042} \\
    DLinear & n/a & \num{13.25} & \num{12.00} & \num{11.50} & \num{12.25} & \num{-0.126} & \num{-0.005} \\
    AutoETS & n/a & \num{11.62} & \num{14.44} & \num{15.81} & \num{13.96} & \num{-0.255} & \num{-0.180} \\
    AutoTheta & n/a & \num{12.94} & \num{15.25} & \num{15.81} & \num{14.67} & \num{-0.268} & \num{-0.188} \\
    \addlinespace
    \multicolumn{8}{l}{\emph{Naive references}} \\
    Climatology & n/a & \num{14.62} & \num{8.50} & \num{7.25} & \num{10.12} & \num{+0.000} & \num{+0.000} \\
    Persistence & n/a & \num{12.38} & \num{16.69} & \num{17.38} & \num{15.48} & \num{-0.444} & \num{-0.278} \\
    Seasonal naive & n/a & \num{17.94} & \num{17.12} & \num{15.25} & \num{16.77} & \num{-0.406} & \num{-0.348} \\
    \bottomrule
  \end{tabular}
\end{table}


\begin{figure}[htbp]
  \centering
  \includegraphics[width=\linewidth]{F6_critical_difference.pdf}
  \caption{Average rank over the 16 daily tasks, with the shaded band marking one Nemenyi critical difference from
  the best model. Foundation models in blue, baselines in grey. The leading group is not internally separable.}
  \label{fig:cd}
\end{figure}

\subsection{Nothing beats climatology on tropical rainfall}

All five baseline wins are a day-of-year climatology on rainfall at 7- and 30-day leads. On the 35-station BMD
record no foundation model exceeds climatology beyond a one-day lead; the best any manages is $-0.004$ CRPSS at
lead 30. On the independent CHIRPS dekadal task none beats it at any horizon. The closest is Chronos-2 given 28
years of context, which matches climatology to within $0.001$ CRPSS (Table~\ref{tab:dekadal}).

The claim needs one scope marker. On the smaller 10-station GHCN-Bangladesh track the best foundation model does
hold positive point estimates at every lead: Toto-2.0 $+0.095$ at day 1, Chronos-2 $+0.019$ at day 30. None of
those margins survives correction ($p = 0.88$, $1.00$, $1.00$), so they are not separable from climatology
either.

\begin{table}[htbp]
  \centering
  \caption{CHIRPS dekadal rainfall, skill against a dekad-of-year climatology (higher is better; 0 is climatology).
  Headline context is 144 dekads; the final column gives the best result with all available history.}
  \label{tab:dekadal}
  \begin{tabular}{lrrrrr}
    \toprule
    Model & Lead 1 & Lead 2 & Lead 3 & Lead 6 & Best, 1{,}024 dekads \\
    \midrule
    Climatology        & \num{0.000} & \num{0.000} & \num{0.000} & \num{0.000} & n/a \\
    Chronos-2          & \num{-0.063} & \num{-0.083} & \num{-0.080} & \num{-0.086} & \num{-0.001} \\
    Toto-2.0           & \num{-0.086} & \num{-0.115} & \num{-0.114} & \num{-0.113} & \num{-0.004} \\
    TimesFM 2.5        & \num{-0.162} & \num{-0.150} & \num{-0.137} & \num{-0.135} & \num{-0.055} \\
    TiRex              & \num{-0.169} & \num{-0.219} & \num{-0.257} & \num{-0.323} & \num{-0.030} \\
    Moirai-2           & \num{-0.196} & \num{-0.290} & \num{-0.347} & \num{-0.504} & \num{-0.073} \\
    Persistence        & \num{-0.478} & \num{-0.635} & \num{-0.763} & \num{-1.158} & n/a \\
    TTM r2 (point only)& \num{-0.682} & \num{-0.868} & \num{-1.002} & \num{-1.529} & \num{-0.846} \\
    \bottomrule
  \end{tabular}
\end{table}

\subsection{Reanalysis flatters models against the stations it represents}

Scoring the same models on NASA POWER and at the co-located stations, skill is significantly higher on reanalysis
in 78 of 108 model-variable-lead comparisons (Figure~\ref{fig:obsrean}). The gap is largest for humidity, significant at every lead for all
nine models tested, and for rainfall, significant for all nine at leads 1 and 30. It is smaller for temperature.
It is weakest for sunshine against radiation, two quantities that are not physically identical, which is why we
do not lean on that row. The effect holds across model families, so it is a property of the target rather
than of any architecture. A benchmark built on reanalysis therefore overstates two things at once: absolute skill, and the
margin foundation models hold over baselines. At a 7-day lead the best foundation model's advantage over the best
trained baseline is $+0.026$ on observations, against $+0.048$ on reanalysis. This is the time-series counterpart of the
critique \citet{gupta2026mausam} make for AI weather-prediction models.

\begin{figure}[htbp]
  \centering
  \includegraphics[width=0.72\linewidth]{F12_obs_vs_reanalysis.pdf}
  \caption{Skill on NASA POWER minus skill at co-located stations, with 95\% moving-block intervals. Positive
  values mean the reanalysis looks more predictable than the observations it represents.}
  \label{fig:obsrean}
\end{figure}

\subsection{A long-lead tropical temperature deficit, mostly explained by sparse context}

Against a trained-baseline control subjected to the same Bangladesh-like missingness, four of the eight foundation
models retain a significant temperature penalty at a 30-day lead: TTM r2 ($+0.187$), Chronos-Bolt ($+0.138$),
Chronos-2 ($+0.060$) and TiRex ($+0.052$). Once both sides face the same gaps, the median penalty across the roster falls from $+0.112$ to
$+0.043$. Roughly \SI{60}{\percent} of the raw effect is therefore attributable to sparser station context rather
than to the tropics (Figure~\ref{fig:did}). At a one-day lead no model shows a
significant penalty.

Two scope markers matter. First, masking only the model, as an ablation applied to the foundation models alone
would, \emph{understates} the effect, because the control also degrades on gappy contexts; the like-for-like
comparison raises the median from $+0.022$ to $+0.043$ and returns TiRex to significance. Second, rainfall shows no
systematic penalty at any lead. Where results are significant at lead 1 the sign is reversed, with four models
performing relatively \emph{better} in Bangladesh than at temperate stations. The single exception is Toto-2.0 at
lead 7 ($+0.028$, CI $0.005$ to $0.053$).

Longer context does not dissolve what remains. Repeating the contrast with 2{,}048 days of history and the control
masked identically leaves the median penalty at $+0.052$, unchanged from 1{,}024 days; Chronos-2 ($+0.059$) and
TiRex ($+0.052$) stay significant, while TimesFM's penalty, which was genuinely context-related, disappears.
An unmasked comparison would have suggested the opposite, because without masking the median halves from $+0.080$
to $+0.043$ as context doubles. That apparent context effect is mostly the missingness effect in disguise.

\begin{figure}[htbp]
  \centering
  \includegraphics[width=0.92\linewidth]{F8_equity_did.pdf}
  \caption{Difference-in-differences against the trained-baseline control, with 95\% intervals. Positive values
  mean the model loses more skill in the tropics than the control does. Masking the control as well as the model is
  the like-for-like comparison.}
  \label{fig:did}
\end{figure}

\subsection{Skill reverses by monsoon regime}

Aggregate leaderboards hide the benchmark's sharpest result. Stratifying July--August targets by spell, and
blocking intervals by spell \emph{event} rather than station-day, 61 of 72 contrasts are significant. In break spells every model is better than climatology, and the foundation models hold the widest
margins: Chronos-Bolt reaches $+0.440$ on temperature at a 30-day lead, against PatchTST's $+0.106$. In active
spells the ordering inverts. Every model falls below climatology at 7- and 30-day leads, and the foundation models
lose the most ground: Chronos-Bolt drops to $-0.989$, against PatchTST's $-0.193$. Humidity repeats the pattern.
Forecast quality is therefore lowest during the very episodes that drive flooding, transplanting and evacuation
decisions.

\subsection{Calibration and compute}

Calibration varies more than accuracy. TimesFM 2.5 (0.831), Toto-2.0 (0.826) and TiRex (0.797) are close to the
nominal \SI{80}{\percent}; Chronos-2, the most accurate model overall, is sharp on rainfall (0.629); and Sundial is
severely under-dispersed (0.323). A GPU diagnostic confirms the last is a property of the model rather than of our
adapter. Interval width is flat in the number of sampled trajectories: 1.44, 1.56 and 1.59 for 20, 100 and 500
samples. The output shape matches the documented layout. And the trajectory spread of \SI{0.64}{\celsius} sits
against the model's own RMSE of \SI{1.51}{\celsius}.

The whole zero-shot leaderboard is cheap to reproduce. Every model fits in \SI{4.2}{\gibi\byte} of GPU memory. All 8 models over
all 16 tasks cost 4.16 GPU-hours on a free Tesla T4, ranging from 0.6 minutes for TTM r2 to 130 minutes for
Sundial. By contrast AutoARIMA alone costs about 20.7 CPU-hours. Accuracy does not track size: TiRex
ranks second overall with 35.3M parameters and \SI{568}{\mebi\byte}, while Toto-2.0 (312.7M) ranks fourth and
Sundial (128.3M) eleventh (Figure~\ref{fig:pareto}).

\begin{figure}[htbp]
  \centering
  \includegraphics[width=0.68\linewidth]{F11_compute_pareto.pdf}
  \caption{Compute against skill on the observation tracks at a 7-day lead; marker area scales with parameter
  count. The full leaderboard is reproducible in roughly four GPU-hours on free hardware.}
  \label{fig:pareto}
\end{figure}
